\section{Three ways to write one estimator}
\label{sec:theory}

\subsection{Phase model}
Following \citet{gomba2016toward} (see \citet{rosen2000synthetic} for interferometry in general), the interferometric phase of band $X$ at centre frequency $f_X$, referenced to the
main band $A$, is
\begin{equation}
  \phi_X = \frac{f_X}{\fA}\,\phi_{nd} + \frac{\fA}{f_X}\,\phi_{io} + 2\pi n_X ,
  \label{eq:model}
\end{equation}
where $\phi_{nd}$ holds everything proportional to frequency (deformation, topography, troposphere), $\phi_{io}$ is the
ionospheric (dispersive) phase at $\fA$, and $n_X$ is the integer an unwrapping assigns. The wrapped phase
is $\psi_X = \W(\phi_X)$, with $\W$ the wrap to $[-\pi,\pi)$. NISAR's 40\,+\,5\,MHz mode has $\fA=1239$\,MHz and
$\fB=1293.5$\,MHz, so $\Delta f/\fA = (\fB-\fA)/\fA = 0.0440$ \citep{nisar2018handbook}.

\subsection{The two-band, reformulated and reordered forms}
Eliminating $\phi_{nd}$ between the two bands gives three forms of the same estimate, which differ only in the order
of subtraction and wrapping.

\paragraph{Two-band.} Both bands are unwrapped and combined,
\begin{equation}
  \hat\phi_{io} = a\,\phiA + b\,\phiB, \qquad a = \frac{\fB^2}{\fB^2-\fA^2} = 12.12,\quad
  b = -\frac{\fA\fB}{\fB^2-\fA^2} = -11.61 .
  \label{eq:twoband}
\end{equation}
This is the classic split-spectrum estimate
\citep{rosen2010measurement,brcic2010estimation,gomba2016toward,fattahi2017insar,liang2019ionospheric} and ISCE3's
\texttt{main\_side\_band} \citep{isce3}.

\paragraph{Reformulated.} \citet{wegmuller2018reformulating} rewrote it in terms of the full-band phase and the
double difference $\Delta = \phiA - \phiB$:
\begin{equation}
  \hat\phi_{io} = c_0\,\phiA + c_1\,\Delta, \qquad c_0 = \frac{\fB}{\fA+\fB} = 0.511,\quad
  c_1 = \frac{\fA\fB}{\fB^2-\fA^2} = 11.61 .
  \label{eq:reform}
\end{equation}
The unwrapped phase is scaled by only $\approx 0.5$, and $\Delta$ is small, so that ``unwrapping can be done to the
lowest ambiguity'' \citep{wegmuller2018reformulating}; \citet{wegmuller2024palsar3} extended it to NISAR's and
PALSAR-3's separate main and side bands. ISCE3's \texttt{main\_diff\_ms\_band}, which makes the ionospheric layer of
the NISAR GUNW product, is this form: it forms $\W(\psiA-\psiB)$, unwraps it with SNAPHU
\citep{chen2001snaphu,chen2002snaphu}, and combines it with the unwrapped main band \citep{isce3}. But $\Delta$
is not free of deformation. From \eqref{eq:model},
\begin{equation}
  \Delta = -\frac{\Delta f}{\fA}\,\phi_{nd} + \frac{\Delta f}{\fB}\,\phi_{io} + 2\pi(n_A - n_B),
  \label{eq:delta}
\end{equation}
so it carries 4.4\% of the deformation, with the deformation's steps and gradients. That content is real enough that \citet{ozawa2022dsi} use the difference of split-band interferograms to measure large coseismic deformation next to the 2016 Kumamoto ruptures, where the full-band phase cannot be unwrapped.

\paragraph{Reordered (ours).} Unwrap the main band only, and remove its scaled phase from the \emph{wrapped} side band
before anything else is unwrapped:
\begin{equation}
  \begin{gathered}
  r = \W\!\left(\psiB - \frac{\fB}{\fA}\,\phiA\right) = \W\!\left(k\,\phi_{io} - 2\pi\,\frac{\Delta f}{\fA}\,n_A\right),\\
  k = \frac{\fA}{\fB} - \frac{\fB}{\fA} = -0.0861, \qquad \hat\phi_{io} = \U(r)/k .
  \end{gathered}
  \label{eq:ours}
\end{equation}
The non-dispersive phase cancels exactly, deformation included, whatever its size or steepness; $\psiB$ enters only
wrapped, so the side band's integer never appears. $r$ is $0.086$ times the ionosphere, smooth and continuous across
the rupture, and it is the only field that is unwrapped. Since $c_0 = c_1\,\Delta f/\fA$ and $c_1 = -1/k$, the
reordered estimate equals \eqref{eq:reform} identically,
\begin{equation}
  \frac{r}{k} = \frac{1}{k}\left(\psiB - \phiA - \frac{\Delta f}{\fA}\phiA\right) = c_0\,\phiA + c_1\,\Delta
  \pmod{2\pi/|k|},
\end{equation}
whenever both unwrappings are right. The three forms are therefore one estimator. They differ in what the unwrapper
is asked to see: $\phiB$ carries all of the deformation, $\Delta$ 4.4\% of it, $r$ none.

\subsection{What an integer error costs}
Table~\ref{tab:integers} lists the raw estimate's response to one wrong cycle in each unwrapped quantity. The
reformulated and reordered forms have identical sensitivities. A cycle of the main band costs $2\pi c_0 =
3.21$\,rad in both, the common-error term $\pi m$ of \citet[their eq.~25]{gomba2016toward}. A cycle of the
second unwrapped field ($\Delta$ or $r$) costs $2\pi/|k| = 73$\,rad in both, or 1.4\,m of LOS. What differs is how
likely that cycle is: $\Delta$ must be unwrapped through the deformation, $r$ need not be. In the two-band form a
main-band cycle that the side band does not share is a \emph{differential} error worth $2\pi a = 76$\,rad.

\begin{table}[t]
\centering
\caption{The raw dispersive estimate's change for one cycle of error in each unwrapped field, at NISAR's
1239/1293.5\,MHz. LOS equivalents are $\lambda_A/4\pi = 19.3$\,mm per radian.}
\label{tab:integers}
\begin{tabular}{lccc}
\toprule
error & two-band \eqref{eq:twoband} & reformulated \eqref{eq:reform} & reordered \eqref{eq:ours} \\
\midrule
main band $n_A$ only            & $2\pi a$ = 76.2\,rad & $2\pi c_0$ = 3.21\,rad & $2\pi c_0$ = 3.21\,rad \\
side band $n_B$ only            & $2\pi b$ = 72.9\,rad & --- (enters wrapped)   & --- (enters wrapped)   \\
both bands, $n_A=n_B$           & $2\pi(a+b)$ = 3.21\,rad & ---                 & ---                    \\
double difference $\Delta$      & ---                  & $2\pi c_1$ = 72.9\,rad & ---                    \\
residual $r$                    & ---                  & ---                    & $2\pi/|k|$ = 72.9\,rad \\
\bottomrule
\end{tabular}
\end{table}

\subsection{Why the double difference is hard to unwrap at a rupture}
Equation~\eqref{eq:delta} puts $-0.044\,\phi_{nd}$ into $\Delta$. A 4\,m strike-slip rupture seen at $\sim$0.6 of
its horizontal motion makes a LOS step of $\sim$2.4\,m, about 125\,rad at L band, which becomes a 5.5\,rad step in
$\Delta$, close to a cycle. Three things then fail.
\begin{enumerate}
  \item Taking $\Delta$ ``to the lowest ambiguity'' requires $|\Delta|<\pi$ about a reference
        \citep{wegmuller2018reformulating}. Its deformation part alone reaches that at
        $|\phi_{nd}| = \pi/0.044 = 71$\,rad, or 1.4\,m of LOS. The authors note that strong ionospheric gradients
        break it too \citep{wegmuller2024palsar3}.
  \item Across a fault zone that is decorrelated, or excluded from the estimate, nothing constrains $\Delta$'s integer.
        An unwrapper that bridges the gap takes the shortest arc, and is a cycle off once the step exceeds $\pi$.
        Where the rupture splits the scene (it runs to a coast or to the swath edge) there is no path around its
        end, and the sides' relative integer is arbitrary.
  \item Close to the trace the deformation's gradient is aliased in the cells. $\Delta$ inherits 4.4\% of it, and
        its own unwrapping errors are each 73\,rad of screen.
\end{enumerate}
None of this touches $r$, because \eqref{eq:ours} removes the deformation at the resolution of the cells before
wrapping. The integer it needs across the rupture is the main band's, which a deformation study must determine
anyway (Section~\ref{sec:tie}).

\subsection{The cycle-correction passes}
\label{sec:corrections}
Split-spectrum processors correct unwrapping errors after the fact. \citet[eqs.~28--29]{gomba2016toward} round the
sub-band difference, and the sum, against the filtered first-pass estimates of $\phi_{nd}$ and $\phi_{io}$. They
note that the common-error correction ``is not robust'': inaccurate estimates ``are misinterpreted as phase unwrapping
errors''. ISCE3, which made the NISAR PR GUNWs with \texttt{unwrap\_correction: True} (the run configuration stored
in each product), applies two passes. First, every unwrapped field (here $\phiA$ and $\Delta$) is rounded against
its own truncated-Gaussian low-pass ($\sigma$ of 4 by 37 side-band pixels, about 2.0 by 2.6\,km). Second, $\phiA$ is
rounded against the first pass's filtered $\hat\phi_{nd}+\hat\phi_{io}$, and $\Delta$ against its prediction from
the two \citep{isce3}. Near a rupture the deformation departs from any 2\,km low-pass by more than $\pi$ wherever its
curvature is large. Each such departure is read as a cycle error and removed, adding $\mp 3.21$\,rad to the raw
estimate (Table~\ref{tab:integers}) with the sign of the deformation's curvature. The result is a scaled,
sign-flipped copy of the high-passed deformation, which is what the GUNW screens show near the 2026 ruptures
(Section~\ref{sec:nisar}). The reordered form needs no such pass: $r$ holds no deformation, so there is nothing for
a correction to mistake.

\subsection{Band co-location}
\label{sec:colocation}
The three forms assume that both bands sample the same ground. If the side band is sampled $\delta$ further down range
than the main band, $\phiB$ picks up $\delta\,(\fB/\fA)\,\partial\phi_{nd}/\partial\rho$, and every form's estimate
picks up
\begin{equation}
  \epsilon_\delta \approx b\,\frac{\fB}{\fA}\,\delta\,\frac{\partial\phi_{nd}}{\partial\rho}
  = -12.1\,\delta\,\frac{\partial\phi_{nd}}{\partial\rho} .
  \label{eq:offset}
\end{equation}
ISCE3 decimates the main band to the side band's range columns by stride. That leaves a constant offset of about 3 main-band pixels ($\delta\approx9$\,m) and footprints 8 times apart (40\,m against 325\,m of slant range). How much that matters depends on the deformation's gradient \emph{down range}. A strike-slip rupture that a polar-orbiting radar sees well strikes roughly east--west, nearly along the range direction, so its across-fault gradient barely enters. In our synthetic geometry (Section~\ref{sec:synthetic}) the range direction lies within 2$^\circ$ of the strike, and a 9\,m offset adds 0.3\,mm of LOS at 2--5\,km from the trace. A rupture whose LOS gradient points down range, such as a thrust striking north--south, would see the full across-fault gradient: 10\,rad/km and $\delta = 9$\,m give 1.1\,rad (21\,mm) by \eqref{eq:offset}. The reordered estimator sums both bands over identical map cells, so $\delta=0$ by construction. What remains is the bands' different resolution inside a cell, which the fault-zone exclusion handles (Section~\ref{sec:estimator}). (Sub-bands cut from one band have a wavenumber shift of their own, which matters for the same reason \citep{li2024wavenumber}.)
